\chapter{1918:Fritz Haber}
\label{chap:chemistry1918}

The Nobel Prize in Chemistry for the year 1918 was awarded to Fritz Haber.

\textbf{Motivation:}
for the synthesis of ammonia from its elements

\section{Fritz Haber}

\subsection*{Early Life and Education}

Fritz Haber was born on 1868-12-09 in Breslau, Prussia (now Wroclaw, Poland).
Fritz Jakob Haber was a German chemist who received the Nobel Prize in Chemistry in 1918 for his invention of the Haber process, a method used in industry to synthesize ammonia from nitrogen gas and hydrogen gas. This invention is important for the large-scale synthesis of fertilizers and explosives. It is estimated that a third of annual global food production uses ammonia from the Haber–Bosch process, and that this food supports nearly half the world's population. For this work, Haber has been called one of the most important scientists and industrial chemists in human history.

\subsection*{Career and Major Research}

Haber was professor at the Technische Hochschule in Karlsruhe. He developed a method for synthesizing ammonia directly from nitrogen and hydrogen at high pressure over a catalyst, first demonstrated in 1909. The process, industrialized with Carl Bosch at BASF, made possible the large-scale production of nitrogen fertilizers and altered the course of agriculture.

\subsection*{Nobel Prize-winning Work}

The work for which Fritz Haber was awarded the Nobel Prize in Chemistry in 1918 was described by the Nobel Committee as: "for the synthesis of ammonia from its elements".

This research fundamentally advanced the field by making groundbreaking contributions that transformed chemical science.

\subsection*{Significance and Impact}

The impact of Fritz Haber's work extends far beyond the specific achievement recognized by the Nobel Prize.
These contributions continue to influence chemical research and education today.

\subsection*{Later Career and Honors}

Fritz Haber continued his scientific work until 1934-01-29, passing away in Basel, Switzerland.
He received numerous honors including membership in major scientific academies and societies.

\subsection*{Legacy}

The legacy of Fritz Haber endures through the lasting impact of his scientific contributions.
Fritz Haber's work continues to be cited and built upon by researchers across the chemical sciences.

\subsection*{Selected Publications}

Fritz Haber's major publications include numerous papers in leading journals of the era, detailing the experimental and theoretical work summarized above.

\clearpage

\section*{Chapter Summary}

Fritz Haber was awarded the prize for the synthesis of ammonia from its elements, making possible the industrial production of nitrogen fertilizers.
